\begin{table}
\caption{Why Curie interval coverage varies between splits, and distance-aware conformal variants}
\label{tab:tableS6_design_effect_distance_conformal}
\tablefont
\begin{tabularx}{\textwidth}{@{}>{\raggedright\arraybackslash}Xrrrrrr@{}}
\toprule
\textbf{Method / variant} & \textbf{Design effect} & \textbf{Effective size} & \textbf{Implied SD} & \textbf{Observed SD} & \textbf{Mean coverage} & \textbf{Farthest decile} \\
\midrule
Split conformal & 13.7 & 305 & 0.5 / 1.7 & 6.0 & -- & -- \\
Normalized conformal & 12.4 & 337 & 0.5 / 1.6 & 2.5 & -- & -- \\
CQR & 16.7 & 251 & 0.5 / 1.9 & 5.8 & -- & -- \\
Normalized, plain & -- & -- & -- & 2.5 & 90.3 & 88.2 \\
Normalized, calibrated per distance quartile & -- & -- & -- & 2.3 & 89.5 & 90.3 \\
Normalized, weighted for covariate shift & -- & -- & -- & 2.4 & 90.3 & 88.2 \\
\bottomrule
\end{tabularx}
\par\vspace{4pt}\raggedright\noindent Design effect from the within-formula intraclass correlation of the calibration exceedance indicator. Effective size = effective number of independent calibration units. Implied SD of training-conditional coverage: for exchangeable rows / for the effective size. Observed SD across 10 splits. SD and coverage in percentage points; farthest decile = coverage in the farthest decile of composition distance.
\end{table}
